% Macros: \R=\mathbb R, \eps=\varepsilon, \one=\mathbf 1 (indicator), \TV=total variation, \dist, \supp, \Lip as usual.
% Section conventions:
All sets in this section lie in fixed bounded regions. Directions are unit
vectors, with antipodal vectors identified when discussing lines. A finite
number of projective charts removes this distinction. For a bounded set
$A$, write $|A|_r$ for its covering number by balls of radius $r$. A tube has
bounded length and the stated perpendicular width; enlarging its length by
a fixed constant is always permitted. We use the ball-bound and subpower
conventions of Section~\ref{sec:conventions}. Bounds for radii in a fixed
larger bounded interval follow by changing constants. In particular,
incidence mass at least subpower means a lower bound $\delta^{o(1)}$.


\begin{theorem}[Linear projection gain]\label{ang:linear-gain}
Fix $k\ge2$ and
\[
 k-1\le q\le k-\tfrac12<t<k.
\]
Fix a tube-width constant $C$. In the case $q=k-1$, also fix
$\zeta>0$. There exist $\epsilon,c>0$ such that, for all sufficiently
small $\delta$, the following holds. Let $\mu$ be a bounded point
subprobability in $\R^{k+1}$ and $\nu$ a direction subprobability on $S^k$
with respective $t$- and $q$-ball bounds of loss $\delta^{-\epsilon}$
down to $\delta$. At $q=k-1$ assume in addition
\begin{equation}\label{ang:endpoint-plate}
 \nu\{w:\dist(w,H)\le r\}\le\delta^{-\epsilon}r^\zeta
       \qquad(\delta\le r\le1)
\end{equation}
for every linear hyperplane $H$. Let $E\subset\R^{k+1}\times S^k$ be a
measurable incidence relation, and write $E_w=\{x:(x,w)\in E\}$.
If each direction fiber $E_w$ is covered by at most
$\delta^{-q-\epsilon}$ tubes parallel to $w$ of width $C\delta$, then
\[
 (\mu\otimes\nu)(E)\le\delta^c.
\]
All constants are uniform over the measures and the incidence relations
with these fixed parameters.
\end{theorem}
